\documentclass[11pt,a4paper]{article}
\usepackage[utf8]{inputenc}
\usepackage[T1]{fontenc}
\usepackage[spanish,es-minimal,es-noshorthands]{babel}
\AtBeginDocument{\renewcommand{\sin}{\operatorname{sen}}\renewcommand{\arcsin}{\operatorname{arcsen}}}
\usepackage{lmodern}
\usepackage[sfdefault,lining]{FiraSans}
\usepackage[a4paper,top=2.3cm,bottom=2.2cm,left=1.8cm,right=1.8cm,headheight=1.3cm,headsep=0.4cm,footskip=1cm]{geometry}
\usepackage{amsmath,amssymb}
\usepackage{icomma}
\usepackage{graphicx}
\usepackage[table]{xcolor}
\usepackage{tikz}
\usepackage{fontawesome5}
\usepackage{ccicons}
\usepackage{fancyhdr}
\usepackage{lastpage}
\usepackage{tasks}
\usepackage{needspace}
\usepackage[most]{tcolorbox}
\usepackage{array}
\usepackage[hidelinks]{hyperref}
\AtBeginDocument{\hypersetup{pdftitle={Ficha de ejercicios de la unidad 10: Estadística. 1º ESO},pdfsubject={Matemáticas 1º ESO},pdfkeywords={estadística, población, muestra, variables estadísticas, frecuencias, tablas de frecuencias, diagrama de barras, polígono de frecuencias, histograma, diagrama de sectores, pictograma, moda, media, mediana, rango, desviación media}}}

\definecolor{navy}{HTML}{1B2748}
\definecolor{accent}{HTML}{E8603A}
\definecolor{teal}{HTML}{1F9E90}
\definecolor{soft}{HTML}{EEF1F8}
\definecolor{muted}{HTML}{5D6685}

\setlength{\parindent}{0pt}
\setlength{\parskip}{3pt}
\settasks{label=\alph*), label-format=\bfseries\color{navy}, label-width=1.3em, item-indent=1.8em, column-sep=0.6em, after-item-skip=5pt, before-skip=3pt, after-skip=4pt}

% ---------- capçalera i peu ----------
\newcommand{\llicencia}{\href{https://creativecommons.org/licenses/by-nc-sa/4.0/deed.es}{\ccbyncsa}\ \ Licencia Creative Commons Atribución-NoComercial-CompartirIgual 4.0 Internacional (CC BY-NC-SA 4.0)}
\pagestyle{fancy}
\fancyhf{}
\fancyhead[L]{\raisebox{-0.35\height}{\includegraphics[height=0.95cm]{logo-gva.png}}}
\fancyhead[C]{\small\color{muted}Ficha de ejercicios · Unidad 10 · 1\textsuperscript{o} ESO}
\fancyhead[R]{\raisebox{-0.35\height}{\includegraphics[height=0.95cm]{logo-vermellar.png}}}
\fancyfoot[L]{\scriptsize\color{muted}\llicencia}
\fancyfoot[R]{\small\color{muted}\thepage}
\renewcommand{\headrulewidth}{0.4pt}
\renewcommand{\headrule}{\color{muted!40}\hrule height \headrulewidth}
\fancypagestyle{primera}{\fancyhf{}\renewcommand{\headrulewidth}{0pt}
  \fancyfoot[L]{\scriptsize\color{muted}\llicencia}
  \fancyfoot[R]{\small\color{muted}\thepage}}

% ---------- nivells i exercicis ----------
\newcommand{\estrelles}[1]{{\color{accent}\ifcase#1\or\faStar\faStar[regular]\faStar[regular]\or\faStar\faStar\faStar[regular]\or\faStar\faStar\faStar\fi}}
\newcounter{exer}
\newcommand{\ex}[1]{\par\needspace{4\baselineskip}\medskip\stepcounter{exer}\noindent{\color{navy}\bfseries\theexer.}\ \ {\small\estrelles{#1}}\ \ }

\newcommand{\seccio}[2]{\par\needspace{10\baselineskip}\bigskip
  \begin{tcolorbox}[colback=navy,colframe=navy,coltext=white,arc=2mm,boxrule=0pt,left=3mm,right=3mm,top=1.2mm,bottom=1.2mm]
  {\large\bfseries #1\quad #2}
  \end{tcolorbox}\vspace{-2pt}}

\newcolumntype{C}[1]{>{\centering\arraybackslash}p{#1}}
\definecolor{violet}{HTML}{7C6CF0}
\definecolor{gold}{HTML}{E0A21B}

% quadre de recordatori (només a la fitxa d'exercicis)
\newcommand{\recorda}[1]{\par\smallskip\begin{tcolorbox}[colback=accent!7,colframe=accent!55,arc=2mm,boxrule=0.5pt,left=2.5mm,right=2.5mm,top=1mm,bottom=1mm]\small{\color{accent}\faLightbulb[regular]}\ \textbf{Recuerda:}\ #1\end{tcolorbox}}
\newcommand{\res}[1]{#1}\newcommand{\rest}[1]{#1}

\def\fmneed{14}
% =====================================================================
% Capçalera FentMates de les fitxes: un color per a cada curs.
% Els colors són els de la pàgina inicial (LEVELS a assets/js/materials.js);
% el logo ve de _eines/logo/pdf/ (python3 _eines/logo/genera.py).
%
% Ús, al final del preàmbul de la unitat:
%   % =====================================================================
% Capçalera FentMates de les fitxes: un color per a cada curs.
% Els colors són els de la pàgina inicial (LEVELS a assets/js/materials.js);
% el logo ve de _eines/logo/pdf/ (python3 _eines/logo/genera.py).
%
% Ús, al final del preàmbul de la unitat:
%   % =====================================================================
% Capçalera FentMates de les fitxes: un color per a cada curs.
% Els colors són els de la pàgina inicial (LEVELS a assets/js/materials.js);
% el logo ve de _eines/logo/pdf/ (python3 _eines/logo/genera.py).
%
% Ús, al final del preàmbul de la unitat:
%   \input{../comu/fentmates.tex}
%   \fmcurs{1eso}
% i, al començament del document:
%   \capcalera{Matemàtiques}{1r ESO}{Fitxa d'exercicis · Unitat 2}{Divisibilitat}
% Opcional abans de \input: \def\fmneed{14} (espai mínim abans de cada secció).
% =====================================================================
\makeatletter
\def\fm@logodir{../../logo/pdf}
\providecommand{\fmneed}{10}
\def\fm@cursnom{}

% clau del curs, color i color del text damunt del color
\newcommand{\fm@curs}[3]{\@namedef{fm@color@#1}{#2}\@namedef{fm@text@#1}{#3}}
\fm@curs{1eso}{E6B84A}{navy}
\fm@curs{2eso}{3B82F6}{white}
\fm@curs{3eso}{FF7A50}{white}
\fm@curs{4eso-a}{A78BFA}{white}
\fm@curs{4eso-b}{6D28D9}{white}
\fm@curs{1batx-mi}{9B1C1C}{white}
\fm@curs{1batx-ccss}{EF6B6B}{white}
\fm@curs{2batx-mii}{1E6B47}{white}
\fm@curs{2batx-ccss}{2BB3A3}{white}

\newcommand{\fmcurs}[1]{%
  \@ifundefined{fm@color@#1}{\PackageError{fentmates}{Curs desconegut: #1}{}}{}%
  \edef\fm@tmp{\noexpand\definecolor{curs}{HTML}{\@nameuse{fm@color@#1}}%
               \noexpand\colorlet{curstext}{\@nameuse{fm@text@#1}}}\fm@tmp
  \colorlet{cursfosc}{curs!72!black}%
  \def\fm@logo{\fm@logodir/vertical-#1.pdf}%
  \def\fm@marca{\fm@logodir/marca-#1.pdf}}

% ---------- capçalera de la primera pàgina ----------
% #1 matèria (petit), #2 curs, #3 tipus de document i unitat, #4 títol
\newcommand{\capcalera}[4]{%
  \gdef\fm@cursnom{#2}%
  \par\noindent
  \begin{tcolorbox}[colback=curs,colframe=curs,boxrule=0pt,arc=2.2mm,width=3.5cm,height=2.15cm,
      valign=center,left=3mm,right=2mm,top=1mm,bottom=1mm,nobeforeafter,box align=center]
    \color{curstext}{\footnotesize #1}\par\vspace{1pt}{\fontsize{19}{21}\selectfont\bfseries #2\par}%
  \end{tcolorbox}\hspace{5mm}%
  \begin{minipage}[c]{\dimexpr\linewidth-3.5cm-5mm-2.9cm\relax}
    {\color{cursfosc}\small\bfseries\MakeUppercase{#3}}\par\vspace{2pt}
    {\color{navy}\fontsize{20}{23}\selectfont\bfseries #4\par}
  \end{minipage}\hfill
  \begin{minipage}[c]{2.6cm}\raggedleft\includegraphics[width=2.6cm]{\fm@logo}\end{minipage}\par
  \vspace{2pt}}

% ---------- barres de secció del color del curs ----------
\renewcommand{\seccio}[2]{\par\needspace{\fmneed\baselineskip}\bigskip
  \begin{tcolorbox}[colback=curs,colframe=curs,coltext=curstext,arc=2mm,boxrule=0pt,left=3mm,right=3mm,top=1.2mm,bottom=1.2mm]
  {\large\bfseries #1\quad #2}
  \end{tcolorbox}\vspace{-2pt}}

% ---------- estrelles de dificultat del color del curs (a Batxillerat no se'n posen) ----------
\@ifundefined{estrelles}{}{%
  \renewcommand{\estrelles}[1]{{\color{curs}\ifcase#1\or\faStar\faStar[regular]\faStar[regular]\or\faStar\faStar\faStar[regular]\or\faStar\faStar\faStar\fi}}}

% ---------- capçalera i peu de les pàgines ----------
% primera pàgina: logos institucionals al peu
\fancypagestyle{primera}{\fancyhf{}\renewcommand{\headrulewidth}{0pt}
  \fancyfoot[L]{\raisebox{-0.32\height}{\includegraphics[height=0.78cm]{logo-gva.png}}\hspace{5mm}%
    \raisebox{-0.32\height}{\includegraphics[height=0.78cm]{logo-vermellar.png}}}
  \fancyfoot[R]{\scriptsize\color{muted}\href{https://creativecommons.org/licenses/by-nc-sa/4.0/deed.ca}{\ccbyncsa}\ \ CC BY-NC-SA 4.0 · fentmates.org\qquad\small\thepage}}
% la resta: el símbol a l'esquerra i el curs a la dreta, amb el títol al centre com abans
\fancyhead[L]{\raisebox{-0.3\height}{\includegraphics[height=0.62cm]{\fm@marca}}}
\fancyhead[R]{\small\bfseries\color{cursfosc}\fm@cursnom}
\renewcommand{\headrule}{\color{curs!45}\hrule height \headrulewidth}
\makeatother

%   \fmcurs{1eso}
% i, al començament del document:
%   \capcalera{Matemàtiques}{1r ESO}{Fitxa d'exercicis · Unitat 2}{Divisibilitat}
% Opcional abans de \input: \def\fmneed{14} (espai mínim abans de cada secció).
% =====================================================================
\makeatletter
\def\fm@logodir{../../logo/pdf}
\providecommand{\fmneed}{10}
\def\fm@cursnom{}

% clau del curs, color i color del text damunt del color
\newcommand{\fm@curs}[3]{\@namedef{fm@color@#1}{#2}\@namedef{fm@text@#1}{#3}}
\fm@curs{1eso}{E6B84A}{navy}
\fm@curs{2eso}{3B82F6}{white}
\fm@curs{3eso}{FF7A50}{white}
\fm@curs{4eso-a}{A78BFA}{white}
\fm@curs{4eso-b}{6D28D9}{white}
\fm@curs{1batx-mi}{9B1C1C}{white}
\fm@curs{1batx-ccss}{EF6B6B}{white}
\fm@curs{2batx-mii}{1E6B47}{white}
\fm@curs{2batx-ccss}{2BB3A3}{white}

\newcommand{\fmcurs}[1]{%
  \@ifundefined{fm@color@#1}{\PackageError{fentmates}{Curs desconegut: #1}{}}{}%
  \edef\fm@tmp{\noexpand\definecolor{curs}{HTML}{\@nameuse{fm@color@#1}}%
               \noexpand\colorlet{curstext}{\@nameuse{fm@text@#1}}}\fm@tmp
  \colorlet{cursfosc}{curs!72!black}%
  \def\fm@logo{\fm@logodir/vertical-#1.pdf}%
  \def\fm@marca{\fm@logodir/marca-#1.pdf}}

% ---------- capçalera de la primera pàgina ----------
% #1 matèria (petit), #2 curs, #3 tipus de document i unitat, #4 títol
\newcommand{\capcalera}[4]{%
  \gdef\fm@cursnom{#2}%
  \par\noindent
  \begin{tcolorbox}[colback=curs,colframe=curs,boxrule=0pt,arc=2.2mm,width=3.5cm,height=2.15cm,
      valign=center,left=3mm,right=2mm,top=1mm,bottom=1mm,nobeforeafter,box align=center]
    \color{curstext}{\footnotesize #1}\par\vspace{1pt}{\fontsize{19}{21}\selectfont\bfseries #2\par}%
  \end{tcolorbox}\hspace{5mm}%
  \begin{minipage}[c]{\dimexpr\linewidth-3.5cm-5mm-2.9cm\relax}
    {\color{cursfosc}\small\bfseries\MakeUppercase{#3}}\par\vspace{2pt}
    {\color{navy}\fontsize{20}{23}\selectfont\bfseries #4\par}
  \end{minipage}\hfill
  \begin{minipage}[c]{2.6cm}\raggedleft\includegraphics[width=2.6cm]{\fm@logo}\end{minipage}\par
  \vspace{2pt}}

% ---------- barres de secció del color del curs ----------
\renewcommand{\seccio}[2]{\par\needspace{\fmneed\baselineskip}\bigskip
  \begin{tcolorbox}[colback=curs,colframe=curs,coltext=curstext,arc=2mm,boxrule=0pt,left=3mm,right=3mm,top=1.2mm,bottom=1.2mm]
  {\large\bfseries #1\quad #2}
  \end{tcolorbox}\vspace{-2pt}}

% ---------- estrelles de dificultat del color del curs (a Batxillerat no se'n posen) ----------
\@ifundefined{estrelles}{}{%
  \renewcommand{\estrelles}[1]{{\color{curs}\ifcase#1\or\faStar\faStar[regular]\faStar[regular]\or\faStar\faStar\faStar[regular]\or\faStar\faStar\faStar\fi}}}

% ---------- capçalera i peu de les pàgines ----------
% primera pàgina: logos institucionals al peu
\fancypagestyle{primera}{\fancyhf{}\renewcommand{\headrulewidth}{0pt}
  \fancyfoot[L]{\raisebox{-0.32\height}{\includegraphics[height=0.78cm]{logo-gva.png}}\hspace{5mm}%
    \raisebox{-0.32\height}{\includegraphics[height=0.78cm]{logo-vermellar.png}}}
  \fancyfoot[R]{\scriptsize\color{muted}\href{https://creativecommons.org/licenses/by-nc-sa/4.0/deed.ca}{\ccbyncsa}\ \ CC BY-NC-SA 4.0 · fentmates.org\qquad\small\thepage}}
% la resta: el símbol a l'esquerra i el curs a la dreta, amb el títol al centre com abans
\fancyhead[L]{\raisebox{-0.3\height}{\includegraphics[height=0.62cm]{\fm@marca}}}
\fancyhead[R]{\small\bfseries\color{cursfosc}\fm@cursnom}
\renewcommand{\headrule}{\color{curs!45}\hrule height \headrulewidth}
\makeatother

%   \fmcurs{1eso}
% i, al començament del document:
%   \capcalera{Matemàtiques}{1r ESO}{Fitxa d'exercicis · Unitat 2}{Divisibilitat}
% Opcional abans de \input: \def\fmneed{14} (espai mínim abans de cada secció).
% =====================================================================
\makeatletter
\def\fm@logodir{../../logo/pdf}
\providecommand{\fmneed}{10}
\def\fm@cursnom{}

% clau del curs, color i color del text damunt del color
\newcommand{\fm@curs}[3]{\@namedef{fm@color@#1}{#2}\@namedef{fm@text@#1}{#3}}
\fm@curs{1eso}{E6B84A}{navy}
\fm@curs{2eso}{3B82F6}{white}
\fm@curs{3eso}{FF7A50}{white}
\fm@curs{4eso-a}{A78BFA}{white}
\fm@curs{4eso-b}{6D28D9}{white}
\fm@curs{1batx-mi}{9B1C1C}{white}
\fm@curs{1batx-ccss}{EF6B6B}{white}
\fm@curs{2batx-mii}{1E6B47}{white}
\fm@curs{2batx-ccss}{2BB3A3}{white}

\newcommand{\fmcurs}[1]{%
  \@ifundefined{fm@color@#1}{\PackageError{fentmates}{Curs desconegut: #1}{}}{}%
  \edef\fm@tmp{\noexpand\definecolor{curs}{HTML}{\@nameuse{fm@color@#1}}%
               \noexpand\colorlet{curstext}{\@nameuse{fm@text@#1}}}\fm@tmp
  \colorlet{cursfosc}{curs!72!black}%
  \def\fm@logo{\fm@logodir/vertical-#1.pdf}%
  \def\fm@marca{\fm@logodir/marca-#1.pdf}}

% ---------- capçalera de la primera pàgina ----------
% #1 matèria (petit), #2 curs, #3 tipus de document i unitat, #4 títol
\newcommand{\capcalera}[4]{%
  \gdef\fm@cursnom{#2}%
  \par\noindent
  \begin{tcolorbox}[colback=curs,colframe=curs,boxrule=0pt,arc=2.2mm,width=3.5cm,height=2.15cm,
      valign=center,left=3mm,right=2mm,top=1mm,bottom=1mm,nobeforeafter,box align=center]
    \color{curstext}{\footnotesize #1}\par\vspace{1pt}{\fontsize{19}{21}\selectfont\bfseries #2\par}%
  \end{tcolorbox}\hspace{5mm}%
  \begin{minipage}[c]{\dimexpr\linewidth-3.5cm-5mm-2.9cm\relax}
    {\color{cursfosc}\small\bfseries\MakeUppercase{#3}}\par\vspace{2pt}
    {\color{navy}\fontsize{20}{23}\selectfont\bfseries #4\par}
  \end{minipage}\hfill
  \begin{minipage}[c]{2.6cm}\raggedleft\includegraphics[width=2.6cm]{\fm@logo}\end{minipage}\par
  \vspace{2pt}}

% ---------- barres de secció del color del curs ----------
\renewcommand{\seccio}[2]{\par\needspace{\fmneed\baselineskip}\bigskip
  \begin{tcolorbox}[colback=curs,colframe=curs,coltext=curstext,arc=2mm,boxrule=0pt,left=3mm,right=3mm,top=1.2mm,bottom=1.2mm]
  {\large\bfseries #1\quad #2}
  \end{tcolorbox}\vspace{-2pt}}

% ---------- estrelles de dificultat del color del curs (a Batxillerat no se'n posen) ----------
\@ifundefined{estrelles}{}{%
  \renewcommand{\estrelles}[1]{{\color{curs}\ifcase#1\or\faStar\faStar[regular]\faStar[regular]\or\faStar\faStar\faStar[regular]\or\faStar\faStar\faStar\fi}}}

% ---------- capçalera i peu de les pàgines ----------
% primera pàgina: logos institucionals al peu
\fancypagestyle{primera}{\fancyhf{}\renewcommand{\headrulewidth}{0pt}
  \fancyfoot[L]{\raisebox{-0.32\height}{\includegraphics[height=0.78cm]{logo-gva.png}}\hspace{5mm}%
    \raisebox{-0.32\height}{\includegraphics[height=0.78cm]{logo-vermellar.png}}}
  \fancyfoot[R]{\scriptsize\color{muted}\href{https://creativecommons.org/licenses/by-nc-sa/4.0/deed.ca}{\ccbyncsa}\ \ CC BY-NC-SA 4.0 · fentmates.org\qquad\small\thepage}}
% la resta: el símbol a l'esquerra i el curs a la dreta, amb el títol al centre com abans
\fancyhead[L]{\raisebox{-0.3\height}{\includegraphics[height=0.62cm]{\fm@marca}}}
\fancyhead[R]{\small\bfseries\color{cursfosc}\fm@cursnom}
\renewcommand{\headrule}{\color{curs!45}\hrule height \headrulewidth}
\makeatother

\fmcurs{1eso}

\begin{document}
\thispagestyle{primera}

\capcalera{Matemáticas}{1\textsuperscript{o} ESO}{Ficha de ejercicios · Unidad 10}{Estadística}
\vspace{10pt}
Nombre y apellidos: \rule{6.6cm}{0.4pt}\hfill Grupo: \rule{1.6cm}{0.4pt}\hfill Fecha: \rule{2.2cm}{0.4pt}
\vspace{8pt}

\begin{tcolorbox}[colback=soft,colframe=soft,arc=2mm,left=3mm,right=3mm,top=1.5mm,bottom=1.5mm]
\small Ejercicios graduados por niveles de dificultad:\quad \estrelles{1}\ básico \qquad \estrelles{2}\ medio \qquad \estrelles{3}\ avanzado.\\[2pt] Resuélvelos en el cuaderno mostrando todos los pasos. Para hacer gráficos, usa papel cuadriculado, regla y transportador, y pon siempre el título y el nombre de cada eje.
\end{tcolorbox}

% =====================================================
\seccio{1}{Población, muestra y variables estadísticas}

\ex{1} Clasifica estas variables en cualitativas, cuantitativas discretas o cuantitativas continuas:\begin{tasks}(3)
  \task número de hijos
  \task color preferido
  \task tiempo en correr 100\,m
  \task páginas de un libro
  \task profesión
  \task temperatura de una ciudad
  \task número de calzado
  \task lengua materna
  \task capacidad de una botella
\end{tasks}\recorda{Si los valores no son números, la variable es \textbf{cualitativa}. Si son números: se \textbf{cuentan} $\to$ \textbf{discreta}; se \textbf{miden} $\to$ \textbf{continua}.}

\ex{1} Para conocer cuántas horas de televisión ven a la semana los 450 alumnos de un instituto, se encuesta a 60 elegidos al azar. Indica la población, la muestra, el tamaño de la muestra, un individuo y la variable estudiada (y de qué tipo es).

\ex{1} En cada caso, di si se estudia toda la población o una muestra:\begin{tasks}(1)
  \task La tutora calcula la nota media de su clase con las notas de todos los alumnos.
  \task Para saber cuánto duran las pilas de una fábrica, se prueban 50.
  \task Una empresa pregunta a 2\,000 personas a quién votarán en las elecciones.
  \task El entrenador mide la altura de todos los jugadores de su equipo.
\end{tasks}

\ex{2} Explica por qué, en estos estudios, hay que trabajar con una muestra y no con toda la población:\begin{tasks}(1)
  \task la calidad de los yogures que hace una fábrica;
  \task la opinión de todos los jóvenes de España sobre las redes sociales;
  \task el tiempo que tardan en fundirse las bombillas de un lote de 10\,000.
\end{tasks}

\ex{2} Razona si estas muestras son representativas:\begin{tasks}(1)
  \task Para saber cuántas horas de deporte hacen los alumnos de un instituto, se pregunta a 30 personas que salen del gimnasio.
  \task Para saber el plato preferido de los alumnos de un instituto, se eligen al azar 10 alumnos de cada curso.
  \task Para saber qué música gusta más a la gente del pueblo, Laura pregunta a sus amigos.
\end{tasks}

\ex{2} Escribe tres valores posibles de cada variable y clasifícala:\begin{tasks}(3)
  \task peso de un perro
  \task nacionalidad
  \task número de primos
\end{tasks}

\ex{3} Quieres saber cuánto tiempo dedican a leer cada semana los 600 alumnos de tu instituto. Describe cómo harías el estudio: cuál es la población, cómo elegirías una muestra representativa de 60 alumnos, qué pregunta harías y de qué tipo sería la variable.


% =====================================================
\seccio{2}{Frecuencias absolutas y relativas}

\ex{1} Estas son las mascotas que tienen 20 alumnos: \ $0,\ \allowbreak 1,\ \allowbreak 2,\ \allowbreak 1,\ \allowbreak 0,\ \allowbreak 3,\ \allowbreak 1,\ \allowbreak 1,\ \allowbreak 2,\ \allowbreak 0,\ \allowbreak 1,\ \allowbreak 4,\ \allowbreak 2,\ \allowbreak 1,\ \allowbreak 0,\ \allowbreak 1,\ \allowbreak 2,\ \allowbreak 1,\ \allowbreak 3,\ \allowbreak 1$. \ Haz el recuento y la tabla con las frecuencias absolutas, relativas y los porcentajes.\recorda{Frecuencia absoluta $f_i$: veces que aparece un valor. Frecuencia relativa: $h_i=\dfrac{f_i}{N}$. Porcentaje: $h_i\cdot100$.}

\ex{1} Estos son los grupos sanguíneos de 25 personas: \ A, O, O, A, B, O, A, AB, O, O, A, B, O, A, O, O, A, B, O, A, AB, O, A, O, A. \ Haz la tabla de frecuencias absolutas, relativas y porcentajes.

\ex{1} Escribe como porcentaje:\begin{tasks}(4)
  \task $0{,}3$
  \task $0{,}05$
  \task $\frac{3}{4}$
  \task $\frac{7}{20}$
\end{tasks}\par Y como frecuencia relativa decimal:\begin{tasks}(4)
  \task $12\,\%$
  \task $60\,\%$
  \task $7{,}5\,\%$
  \task $100\,\%$
\end{tasks}

\ex{2} Completa la tabla del deporte preferido de 40 alumnos:\par\smallskip\begin{center}{\renewcommand{\arraystretch}{1.25}\begin{tabular}{|c|C{1.2cm}|C{1.2cm}|C{1.2cm}|}\hline Deporte & $f_i$ & $h_i$ & \%\\\hline Fútbol & 14 & & \\\hline Baloncesto & & $0{,}25$ & \\\hline Tenis & & & 15\\\hline Natación & 8 & & \\\hline Otros & & & \\\hline \textbf{Total} & 40 & & \\\hline\end{tabular}}\end{center}

\ex{2} Hemos lanzado un dado 60 veces. Ha salido el 1 nueve veces, el 2 once, el 3 diez, el 4 ocho, el 5 doce y el 6 el resto. ¿Cuántas veces ha salido el 6? ¿En qué porcentaje de los lanzamientos ha salido un número par? ¿Y un número mayor que 4?

\ex{2} En un instituto de 120 alumnos, el $35\,\%$ va a pie. ¿Cuántos alumnos son? Si 18 alumnos van en bici, ¿cuál es la frecuencia relativa de este valor?

\ex{3} En una clase, la frecuencia relativa de los alumnos que llevan gafas es $0{,}2$, y son 6. ¿Cuántos alumnos hay en la clase? ¿Cuántos no llevan gafas?

\ex{3} Una variable toma tres valores, $A$, $B$ y $C$. La frecuencia relativa de $A$ es $0{,}25$; la de $B$, $0{,}4$, y la frecuencia absoluta de $B$ es 16. Calcula $N$, la frecuencia relativa de $C$ y todas las frecuencias absolutas.


% =====================================================
\seccio{3}{Tablas de frecuencias}

\ex{1} Estas son las notas de un examen de 30 alumnos: \ $5,\ \allowbreak 7,\ \allowbreak 6,\ \allowbreak 4,\ \allowbreak 6,\ \allowbreak 8,\ \allowbreak 5,\ \allowbreak 6,\ \allowbreak 3,\ \allowbreak 7,\ \allowbreak 5,\ \allowbreak 9,\ \allowbreak 6,\ \allowbreak 2,\ \allowbreak 7,\ \allowbreak 6,\ \allowbreak 5,\ \allowbreak 8,\ \allowbreak 4,\ \allowbreak 6,\ \allowbreak 10,\ \allowbreak 7,\ \allowbreak 5,\ \allowbreak 3,\ \allowbreak 6,\ \allowbreak 8,\ \allowbreak 4,\ \allowbreak 7,\ \allowbreak 5,\ \allowbreak 9$. \ Haz la tabla de frecuencias con $f_i$, $F_i$, $h_i$ y porcentajes.\recorda{La frecuencia absoluta acumulada $F_i$ es el número de datos menores o iguales que $x_i$: se suman las $f_i$ hasta esa fila. La última es $N$.}

\ex{1} Con la tabla del ejercicio anterior, responde:\begin{tasks}(2)
  \task ¿Cuántos alumnos han sacado un 6 o menos?
  \task ¿Cuántos han aprobado (5 o más)?
  \task ¿Cuántos han sacado más de un 7?
  \task ¿Qué porcentaje ha sacado un 9 o más?
\end{tasks}

\ex{2} En esta tabla, $x_i$ es el número de libros leídos en un mes y $F_i$, la frecuencia acumulada. Halla las frecuencias absolutas. ¿Cuántos alumnos han leído exactamente 3 libros? ¿Y al menos 2?\par\smallskip\begin{center}{\renewcommand{\arraystretch}{1.3}\begin{tabular}{|c|C{0.9cm}|C{0.9cm}|C{0.9cm}|C{0.9cm}|C{0.9cm}|C{0.9cm}|}\hline $x_i$ & $0$ & $1$ & $2$ & $3$ & $4$ & $5$\\\hline $F_i$ & $2$ & $5$ & $9$ & $15$ & $18$ & $20$\\\hline\end{tabular}}\end{center}

\ex{2} Estos son los pesos (en kg) de 24 alumnos: \ $42,\ \allowbreak 51,\ \allowbreak 38,\ \allowbreak 47,\ \allowbreak 55,\ \allowbreak 44,\ \allowbreak 49,\ \allowbreak 36,\ \allowbreak 52,\ \allowbreak 46,\ \allowbreak 41,\ \allowbreak 58,\ \allowbreak 48,\ \allowbreak 45,\ \allowbreak 50,\ \allowbreak 39,\ \allowbreak 53,\ \allowbreak 47,\ \allowbreak 43,\ \allowbreak 56,\ \allowbreak 49,\ \allowbreak 40,\ \allowbreak 45,\ \allowbreak 51$. \ Agrúpalos en intervalos de amplitud 5 desde 35\,kg y haz la tabla con las marcas de clase, $f_i$, $F_i$ y porcentajes.\recorda{El intervalo $[40,\ 45)$ contiene los datos mayores o iguales que 40 y menores que 45. La marca de clase es el punto medio: $\frac{40+45}{2}=42{,}5$.}

\ex{2} ¿En qué intervalo del ejercicio anterior pondrías cada uno de estos datos?\begin{tasks}(4)
  \task $45$\,kg
  \task $44{,}9$\,kg
  \task $55$\,kg
  \task $35$\,kg
\end{tasks}

\ex{3} Estos son los tiempos (en segundos) de 20 alumnos en una carrera de 100\,m: \ $13{,}8,\ \allowbreak 15{,}2,\ \allowbreak 14{,}1,\ \allowbreak 12{,}6,\ \allowbreak 16{,}3,\ \allowbreak 14{,}9,\ \allowbreak 13{,}2,\ \allowbreak 15{,}7,\ \allowbreak 14{,}4,\ \allowbreak 13{,}9,\ \allowbreak 15{,}1,\ \allowbreak 12{,}9,\ \allowbreak 14{,}6,\ \allowbreak 16{,}8,\ \allowbreak 13{,}5,\ \allowbreak 14{,}2,\ \allowbreak 15{,}5,\ \allowbreak 14{,}8,\ \allowbreak 13{,}1,\ \allowbreak 15{,}9$. \ Agrúpalos en intervalos de 1 segundo desde 12\,s y haz la tabla con $f_i$ y $F_i$. ¿Qué porcentaje ha tardado menos de 14\,s? ¿Cuántos han tardado 15\,s o más?

\ex{3} En una tabla de frecuencias, $F_3=12$, $F_4=19$ y $F_5=25$, y 5 es el valor más grande. ¿Cuántos datos hay? ¿Cuál es la frecuencia absoluta del 4? ¿Qué porcentaje de los datos es mayor que 3?


% =====================================================
\seccio{4}{Diagramas de barras, polígonos e histogramas}

\ex{1} Dibuja el diagrama de barras de los grupos sanguíneos del ejercicio 9 (A: 9, O: 11, B: 3, AB: 2).\recorda{En un \textbf{diagrama de barras} ponemos los valores en el eje horizontal y, encima de cada uno, una barra de altura igual a su frecuencia. Las barras van separadas.}

\ex{1} Este diagrama muestra cuántos viajes en avión han hecho los alumnos de una clase. Haz la tabla de frecuencias. ¿Cuántos alumnos hay? ¿Cuál es la moda? ¿Cuántos han viajado en avión alguna vez?\par\smallskip\begin{center}\begin{tikzpicture}[baseline=(current bounding box.center),line cap=round,line join=round]
\draw[muted!30,line width=0.3pt] (0,0.000)--(5.600,0.000); \node[left,font=\tiny,text=muted,inner sep=1.5pt] at (0,0.000) {$0$};
\draw[muted!30,line width=0.3pt] (0,0.350)--(5.600,0.350); \node[left,font=\tiny,text=muted,inner sep=1.5pt] at (0,0.350) {$1$};
\draw[muted!30,line width=0.3pt] (0,0.700)--(5.600,0.700); \node[left,font=\tiny,text=muted,inner sep=1.5pt] at (0,0.700) {$2$};
\draw[muted!30,line width=0.3pt] (0,1.050)--(5.600,1.050); \node[left,font=\tiny,text=muted,inner sep=1.5pt] at (0,1.050) {$3$};
\draw[muted!30,line width=0.3pt] (0,1.400)--(5.600,1.400); \node[left,font=\tiny,text=muted,inner sep=1.5pt] at (0,1.400) {$4$};
\draw[muted!30,line width=0.3pt] (0,1.750)--(5.600,1.750); \node[left,font=\tiny,text=muted,inner sep=1.5pt] at (0,1.750) {$5$};
\draw[muted!30,line width=0.3pt] (0,2.100)--(5.600,2.100); \node[left,font=\tiny,text=muted,inner sep=1.5pt] at (0,2.100) {$6$};
\draw[muted!30,line width=0.3pt] (0,2.450)--(5.600,2.450); \node[left,font=\tiny,text=muted,inner sep=1.5pt] at (0,2.450) {$7$};
\draw[muted!30,line width=0.3pt] (0,2.800)--(5.600,2.800); \node[left,font=\tiny,text=muted,inner sep=1.5pt] at (0,2.800) {$8$};
\draw[-latex,muted,line width=0.6pt] (0,0)--(5.900,0) node[right,font=\scriptsize,text=navy]{Viajes};
\draw[-latex,muted,line width=0.6pt] (0,0)--(0,3.150) node[above,font=\scriptsize,text=navy]{Alumnos};
\draw[teal,fill=teal!35] (0.280,0) rectangle (0.820,1.050);
\node[below,font=\scriptsize,inner sep=2pt] at (0.550,0) {0};
\draw[teal,fill=teal!35] (1.180,0) rectangle (1.720,2.800);
\node[below,font=\scriptsize,inner sep=2pt] at (1.450,0) {1};
\draw[teal,fill=teal!35] (2.080,0) rectangle (2.620,2.100);
\node[below,font=\scriptsize,inner sep=2pt] at (2.350,0) {2};
\draw[teal,fill=teal!35] (2.980,0) rectangle (3.520,1.400);
\node[below,font=\scriptsize,inner sep=2pt] at (3.250,0) {3};
\draw[teal,fill=teal!35] (3.880,0) rectangle (4.420,0.700);
\node[below,font=\scriptsize,inner sep=2pt] at (4.150,0) {4};
\draw[teal,fill=teal!35] (4.780,0) rectangle (5.320,0.350);
\node[below,font=\scriptsize,inner sep=2pt] at (5.050,0) {5};
\end{tikzpicture}\end{center}

\ex{1} Dibuja el diagrama de barras y el polígono de frecuencias del número de mascotas del ejercicio 8.

\ex{2} Dibuja el histograma y el polígono de frecuencias de los pesos agrupados en intervalos del ejercicio 19.\recorda{En un \textbf{histograma} las barras ocupan todo el intervalo y van juntas. El polígono une los puntos de arriba sobre las marcas de clase y se cierra en el eje.}

\ex{2} ¿Qué gráfico elegirías para representar cada variable?\begin{tasks}(1)
  \task el número de hermanos de 30 alumnos;
  \task la altura de 200 personas;
  \task el medio de transporte preferido para ir de vacaciones;
  \task los minutos que tardan 50 alumnos en hacer los deberes.
\end{tasks}

\ex{2} Las notas de un examen de dos clases de 30 alumnos están en la tabla. Dibuja los dos polígonos de frecuencias en el mismo gráfico. ¿Qué clase ha sacado mejores notas?\par\smallskip\begin{center}{\renewcommand{\arraystretch}{1.25}\begin{tabular}{|c|c|c|c|c|c|c|c|}\hline Nota & 4 & 5 & 6 & 7 & 8 & 9 & 10\\\hline Clase A & 1 & 2 & 5 & 7 & 8 & 4 & 3\\\hline Clase B & 0 & 1 & 3 & 6 & 9 & 7 & 4\\\hline\end{tabular}}\end{center}

\ex{3} Tres candidatos de unas elecciones han obtenido 96, 92 y 98 votos. Un cartel los representa con un diagrama de barras en el que el eje vertical empieza en 90. ¿Por qué es engañoso? Dibújalo bien.


% =====================================================
\seccio{5}{Diagramas de sectores y pictogramas}

\ex{1} Se ha preguntado a 36 alumnos por su deporte preferido: fútbol 15, baloncesto 9, tenis 6, natación 4 y otros 2. Calcula el ángulo de cada sector y dibuja el diagrama de sectores.\recorda{Ángulo del sector $=\dfrac{f_i}{N}\cdot360^\circ$. Si $N$ es divisor de 360, cada dato vale $360^\circ:N$.}

\ex{1} Calcula los ángulos del diagrama de sectores de los grupos sanguíneos del ejercicio 9 (A: 9, O: 11, B: 3, AB: 2) y dibújalo.

\ex{2} Se ha preguntado a 200 jóvenes qué aparato usan más para jugar. ¿Cuántas personas han elegido cada aparato?\par\smallskip\begin{center}\begin{tikzpicture}[baseline=(current bounding box.center)]
\fill[teal!70,draw=white,line width=0.8pt] (0,0)--(90.000:1.45) arc[start angle=90.000,end angle=-54.000,radius=1.45cm]--cycle;
\node[font=\scriptsize\bfseries,text=white] at (18.000:0.899) {$40\,\%$};
\fill[accent!70,draw=white,line width=0.8pt] (0,0)--(-54.000:1.45) arc[start angle=-54.000,end angle=-144.000,radius=1.45cm]--cycle;
\node[font=\scriptsize\bfseries,text=white] at (-99.000:0.899) {$25\,\%$};
\fill[violet!70,draw=white,line width=0.8pt] (0,0)--(-144.000:1.45) arc[start angle=-144.000,end angle=-216.000,radius=1.45cm]--cycle;
\node[font=\scriptsize\bfseries,text=white] at (-180.000:0.899) {$20\,\%$};
\fill[gold!70,draw=white,line width=0.8pt] (0,0)--(-216.000:1.45) arc[start angle=-216.000,end angle=-270.000,radius=1.45cm]--cycle;
\node[font=\scriptsize\bfseries,text=white] at (-243.000:0.899) {$15\,\%$};
\fill[teal!70] (1.85,0.415) rectangle ++(0.22,0.22); \node[right,font=\scriptsize] at (2.10,0.525) {Móvil};
\fill[accent!70] (1.85,0.065) rectangle ++(0.22,0.22); \node[right,font=\scriptsize] at (2.10,0.175) {Consola};
\fill[violet!70] (1.85,-0.285) rectangle ++(0.22,0.22); \node[right,font=\scriptsize] at (2.10,-0.175) {Ordenador};
\fill[gold!70] (1.85,-0.635) rectangle ++(0.22,0.22); \node[right,font=\scriptsize] at (2.10,-0.525) {Tableta};
\end{tikzpicture}\end{center}

\ex{2} En un diagrama de sectores sobre el lugar de vacaciones preferido de 72 personas, los ángulos son: playa $100^\circ$, montaña $120^\circ$, ciudad $80^\circ$ y pueblo $60^\circ$. ¿Cuántas personas han elegido cada lugar?\recorda{Del sector a la frecuencia: $f_i=\dfrac{\text{ángulo}}{360^\circ}\cdot N$.}

\ex{2} Este pictograma muestra las bicis que ha vendido una tienda cada mes. Cada icono representa 4 bicis. ¿Cuántas ha vendido cada mes? ¿Y en total?\par\smallskip\begin{center}\begin{tikzpicture}[baseline=(current bounding box.center)]
\node[left,font=\small] at (0,-0.00) {Enero};
\node[text=accent,font=\large] at (0.25,-0.00) {\faBicycle};
\node[text=accent,font=\large] at (0.85,-0.00) {\faBicycle};
\node[text=accent,font=\large] at (1.45,-0.00) {\faBicycle};
\node[left,font=\small] at (0,-0.55) {Febrero};
\node[text=accent,font=\large] at (0.25,-0.55) {\faBicycle};
\node[text=accent,font=\large] at (0.85,-0.55) {\faBicycle};
\begin{scope}\clip (1.15,-0.85) rectangle (1.450,-0.25); \node[text=accent,font=\large] at (1.45,-0.55) {\faBicycle};\end{scope}
\node[left,font=\small] at (0,-1.10) {Marzo};
\node[text=accent,font=\large] at (0.25,-1.10) {\faBicycle};
\node[text=accent,font=\large] at (0.85,-1.10) {\faBicycle};
\node[text=accent,font=\large] at (1.45,-1.10) {\faBicycle};
\node[text=accent,font=\large] at (2.05,-1.10) {\faBicycle};
\node[left,font=\small] at (0,-1.65) {Abril};
\node[text=accent,font=\large] at (0.25,-1.65) {\faBicycle};
\node[text=accent,font=\large] at (0.85,-1.65) {\faBicycle};
\node[text=accent,font=\large] at (1.45,-1.65) {\faBicycle};
\node[text=accent,font=\large] at (2.05,-1.65) {\faBicycle};
\node[text=accent,font=\large] at (2.65,-1.65) {\faBicycle};
\begin{scope}\clip (2.95,-1.95) rectangle (3.250,-1.35); \node[text=accent,font=\large] at (3.25,-1.65) {\faBicycle};\end{scope}
\end{tikzpicture}\end{center}

\ex{2} Los alumnos de cuatro grupos de 1\textsuperscript{o} de ESO han recogido tapones para una campaña solidaria: 1\textsuperscript{o} A, 20 bolsas; 1\textsuperscript{o} B, 35; 1\textsuperscript{o} C, 15, y 1\textsuperscript{o} D, 30. Haz un pictograma en el que cada icono valga 5 bolsas.

\ex{3} En un diagrama de sectores, el sector de las personas que prefieren el verano mide $72^\circ$ y corresponde a 18 personas. ¿Cuántas personas se han encuestado? Si el $30\,\%$ prefiere la primavera, ¿cuántas personas son y qué ángulo tiene su sector?


% =====================================================
\seccio{6}{Moda, media y mediana}

\ex{1} Calcula la moda, la media y la mediana de: \ $5,\ \allowbreak 8,\ \allowbreak 3,\ \allowbreak 8,\ \allowbreak 6,\ \allowbreak 9,\ \allowbreak 7$.\recorda{\textbf{Moda}: el valor más repetido. \textbf{Media}: $\bar x=\dfrac{\text{suma de los datos}}{N}$. \textbf{Mediana}: el valor central de los datos \textbf{ordenados}; si hay dos, su media.}

\ex{1} Calcula la moda, la media y la mediana de: \ $12,\ \allowbreak 15,\ \allowbreak 9,\ \allowbreak 18,\ \allowbreak 15,\ \allowbreak 11$.

\ex{1} Se ha preguntado a 30 alumnos por su estación preferida: primavera 9, verano 13, otoño 3 e invierno 5. ¿Cuál es la moda? ¿Se pueden calcular la media y la mediana?

\ex{2} Calcula la moda, la media y la mediana de cada conjunto de datos:\begin{tasks}(1)
  \task $4,\ \allowbreak 9,\ \allowbreak 6,\ \allowbreak 9,\ \allowbreak 2$
  \task $1,\ \allowbreak 3,\ \allowbreak 3,\ \allowbreak 7,\ \allowbreak 8,\ \allowbreak 8$
  \task $10,\ \allowbreak 20,\ \allowbreak 30,\ \allowbreak 40$
\end{tasks}

\ex{2} Los goles que ha marcado un equipo en los últimos 8 partidos son: \ $2,\ \allowbreak 1,\ \allowbreak 0,\ \allowbreak 3,\ \allowbreak 2,\ \allowbreak 2,\ \allowbreak 1,\ \allowbreak 4$. \ Calcula la moda, la media y la mediana.

\ex{2} Los precios (en miles de euros) de las casas vendidas en una calle han sido: 150, 160, 155, 170, 165 y 900. Calcula la media y la mediana. ¿Cuál de los dos parámetros representa mejor el precio de una casa de la calle? ¿Por qué?

\ex{3} Escribe cinco números que tengan media 6, mediana 5 y moda 4.

\ex{3} Las notas de Juan son 4, 6, 7 y 9. ¿Cuál es la media? ¿Qué nota tiene que sacar en el siguiente examen para que la media sea 7?


% =====================================================
\seccio{7}{Parámetros a partir de una tabla de frecuencias}

\ex{1} El número de hermanos de 25 alumnos es el de la tabla. Calcula la media, la mediana y la moda.\par\smallskip\begin{center}{\renewcommand{\arraystretch}{1.3}\begin{tabular}{|c|C{0.9cm}|C{0.9cm}|C{0.9cm}|C{0.9cm}|C{0.9cm}|}\hline $x_i$ & $0$ & $1$ & $2$ & $3$ & $4$\\\hline $f_i$ & $4$ & $10$ & $6$ & $3$ & $2$\\\hline\end{tabular}}\end{center}\recorda{$\bar x=\dfrac{\sum x_i\cdot f_i}{N}$. Para la mediana, busca la posición central y mira en qué fila la $F_i$ llega a ella.}

\ex{1} Calcula la media, la mediana y la moda de las notas del ejercicio 16 a partir de su tabla de frecuencias.

\ex{2} Esta tabla recoge el número de calzado de 28 alumnos. Calcula la media, la mediana y la moda.\par\smallskip\begin{center}{\renewcommand{\arraystretch}{1.3}\begin{tabular}{|c|C{0.9cm}|C{0.9cm}|C{0.9cm}|C{0.9cm}|C{0.9cm}|C{0.9cm}|}\hline Número & $35$ & $36$ & $37$ & $38$ & $39$ & $40$\\\hline $f_i$ & $2$ & $5$ & $7$ & $8$ & $4$ & $2$\\\hline\end{tabular}}\end{center}

\ex{2} Los goles marcados por un equipo en 30 partidos han sido: 0 goles en 5 partidos, 1 gol en 8, 2 en 9, 3 en 6 y 4 en 2. Calcula la media de goles por partido, la mediana y la moda.

\ex{2} Calcula la media aproximada de los pesos agrupados del ejercicio 19 usando las marcas de clase.

\ex{3} En esta tabla falta una frecuencia. Hállala sabiendo que la media es $1{,}5$.\par\smallskip\begin{center}{\renewcommand{\arraystretch}{1.3}\begin{tabular}{|c|C{0.9cm}|C{0.9cm}|C{0.9cm}|C{0.9cm}|}\hline $x_i$ & $0$ & $1$ & $2$ & $3$\\\hline $f_i$ & $4$ & $a$ & $6$ & $4$\\\hline\end{tabular}}\end{center}

\ex{3} Este diagrama de barras muestra la nota (de 1 a 5) que 20 clientes han puesto a un restaurante. Calcula la media, la mediana y la moda.\par\smallskip\begin{center}\begin{tikzpicture}[baseline=(current bounding box.center),line cap=round,line join=round]
\draw[muted!30,line width=0.3pt] (0,0.000)--(4.700,0.000); \node[left,font=\tiny,text=muted,inner sep=1.5pt] at (0,0.000) {$0$};
\draw[muted!30,line width=0.3pt] (0,0.350)--(4.700,0.350); \node[left,font=\tiny,text=muted,inner sep=1.5pt] at (0,0.350) {$1$};
\draw[muted!30,line width=0.3pt] (0,0.700)--(4.700,0.700); \node[left,font=\tiny,text=muted,inner sep=1.5pt] at (0,0.700) {$2$};
\draw[muted!30,line width=0.3pt] (0,1.050)--(4.700,1.050); \node[left,font=\tiny,text=muted,inner sep=1.5pt] at (0,1.050) {$3$};
\draw[muted!30,line width=0.3pt] (0,1.400)--(4.700,1.400); \node[left,font=\tiny,text=muted,inner sep=1.5pt] at (0,1.400) {$4$};
\draw[muted!30,line width=0.3pt] (0,1.750)--(4.700,1.750); \node[left,font=\tiny,text=muted,inner sep=1.5pt] at (0,1.750) {$5$};
\draw[muted!30,line width=0.3pt] (0,2.100)--(4.700,2.100); \node[left,font=\tiny,text=muted,inner sep=1.5pt] at (0,2.100) {$6$};
\draw[muted!30,line width=0.3pt] (0,2.450)--(4.700,2.450); \node[left,font=\tiny,text=muted,inner sep=1.5pt] at (0,2.450) {$7$};
\draw[muted!30,line width=0.3pt] (0,2.800)--(4.700,2.800); \node[left,font=\tiny,text=muted,inner sep=1.5pt] at (0,2.800) {$8$};
\draw[-latex,muted,line width=0.6pt] (0,0)--(5.000,0) node[right,font=\scriptsize,text=navy]{Nota};
\draw[-latex,muted,line width=0.6pt] (0,0)--(0,3.150) node[above,font=\scriptsize,text=navy]{Clientes};
\draw[teal,fill=teal!35] (0.280,0) rectangle (0.820,0.700);
\node[below,font=\scriptsize,inner sep=2pt] at (0.550,0) {1};
\draw[teal,fill=teal!35] (1.180,0) rectangle (1.720,1.750);
\node[below,font=\scriptsize,inner sep=2pt] at (1.450,0) {2};
\draw[teal,fill=teal!35] (2.080,0) rectangle (2.620,2.800);
\node[below,font=\scriptsize,inner sep=2pt] at (2.350,0) {3};
\draw[teal,fill=teal!35] (2.980,0) rectangle (3.520,1.050);
\node[below,font=\scriptsize,inner sep=2pt] at (3.250,0) {4};
\draw[teal,fill=teal!35] (3.880,0) rectangle (4.420,0.700);
\node[below,font=\scriptsize,inner sep=2pt] at (4.150,0) {5};
\end{tikzpicture}\end{center}


% =====================================================
\seccio{8}{Rango y desviación media}

\ex{1} Calcula el rango de cada conjunto de datos:\begin{tasks}(3)
  \task $3,\ \allowbreak 12,\ \allowbreak 7,\ \allowbreak 15,\ \allowbreak 9$
  \task $-4,\ \allowbreak 6,\ \allowbreak 2,\ \allowbreak -1,\ \allowbreak 8$
  \task $2{,}5,\ \allowbreak 4{,}8,\ \allowbreak 1{,}3,\ \allowbreak 3{,}6$
\end{tasks}\recorda{Rango $=x_{\max}-x_{\min}$. \ Desviación media: $DM=\dfrac{\sum|x_i-\bar x|}{N}$ (la media de las distancias a la media).}

\ex{1} Calcula la desviación media de: \ $3,\ \allowbreak 5,\ \allowbreak 7,\ \allowbreak 9,\ \allowbreak 11$.

\ex{2} Calcula la media, el rango y la desviación media de: \ $2,\ \allowbreak 4,\ \allowbreak 4,\ \allowbreak 6,\ \allowbreak 8,\ \allowbreak 12$.

\ex{2} Las notas de dos alumnos en seis exámenes han sido: Ana, $5,\ \allowbreak 6,\ \allowbreak 5,\ \allowbreak 6,\ \allowbreak 5,\ \allowbreak 6$; Gabriel, $2,\ \allowbreak 9,\ \allowbreak 3,\ \allowbreak 8,\ \allowbreak 4,\ \allowbreak 7$. Calcula la media, el rango y la desviación media de cada uno. ¿Quién es más regular?

\ex{2} Las temperaturas máximas de una semana de cinco días en dos ciudades han sido: ciudad X, $12,\ \allowbreak 14,\ \allowbreak 15,\ \allowbreak 13,\ \allowbreak 16$\,°C; ciudad Y, $6,\ \allowbreak 20,\ \allowbreak 10,\ \allowbreak 18,\ \allowbreak 16$\,°C. Compáralas con la media, el rango y la desviación media.

\ex{3} Calcula la media y la desviación media de los datos de esta tabla.\par\smallskip\begin{center}{\renewcommand{\arraystretch}{1.3}\begin{tabular}{|c|C{0.9cm}|C{0.9cm}|C{0.9cm}|C{0.9cm}|C{0.9cm}|}\hline $x_i$ & $1$ & $2$ & $3$ & $4$ & $5$\\\hline $f_i$ & $2$ & $3$ & $5$ & $3$ & $2$\\\hline\end{tabular}}\end{center}

\ex{3} Razona si es verdadero o falso:\begin{tasks}(1)
  \task Si el rango es 0, todos los datos son iguales.
  \task Si sumamos 5 a todos los datos, la desviación media aumenta 5.
  \task Dos conjuntos con la misma media tienen la misma desviación media.
  \task La desviación media nunca puede ser negativa.
\end{tasks}


% =====================================================
\seccio{9}{Problemas de estadística}

\ex{1} Marina ha sacado $6{,}5$; $7$ y $5{,}5$ en los tres primeros exámenes. ¿Qué nota necesita en el cuarto para tener una media de $6{,}5$?\recorda{Si conocemos la media y el número de datos, conocemos la suma: \ $\text{suma}=\bar x\cdot N$.}

\ex{1} La media de cuatro números es 12. Tres de ellos son 10, 15 y 9. ¿Cuál es el cuarto?

\ex{2} En una clase de 25 alumnos la nota media de un examen ha sido $6{,}4$ y en otra de 15 alumnos, $7{,}2$. ¿Cuál es la nota media de los 40 alumnos? ¿Por qué no es $6{,}8$?

\ex{2} Se ha preguntado a 20 familias cuántos hijos tienen: \ $0,\ \allowbreak 1,\ \allowbreak 2,\ \allowbreak 2,\ \allowbreak 1,\ \allowbreak 3,\ \allowbreak 2,\ \allowbreak 1,\ \allowbreak 1,\ \allowbreak 2,\ \allowbreak 0,\ \allowbreak 2,\ \allowbreak 3,\ \allowbreak 1,\ \allowbreak 2,\ \allowbreak 4,\ \allowbreak 2,\ \allowbreak 1,\ \allowbreak 2,\ \allowbreak 1$.\begin{tasks}(1)
  \task Haz la tabla de frecuencias (con $f_i$, $F_i$ y $h_i$).
  \task Dibuja su diagrama de barras.
  \task Calcula la moda, la media, la mediana y el rango.
  \task ¿Qué porcentaje de familias tiene más de un hijo? ¿Qué ángulo tendría el sector «2 hijos» en un diagrama de sectores?
\end{tasks}

\ex{2} La edad media de los 11 jugadores de un equipo es de 22 años. Si añadimos al entrenador, la edad media de los 12 es de 24 años. ¿Cuántos años tiene el entrenador?

\ex{3} Un profesor ha calculado que la media de las notas de los 10 alumnos de un grupo es $6{,}3$, pero después se da cuenta de que ha copiado un 9 en lugar de un 4. ¿Cuál es la media correcta? ¿Cambia la mediana si las notas eran 3, 4, 5, 5, 6, 7, 7, 8, 9, 9 (con el 9 equivocado)?

\ex{3} De dos clases se conocen estos resultados de un examen: clase A, $\bar x=6$, $Me=6$, $DM=0{,}8$; clase B, $\bar x=6$, $Me=7$, $DM=2{,}5$. Razona si es verdadero o falso:\begin{tasks}(1)
  \task Las dos clases han sacado, de media, la misma nota.
  \task En la clase B, al menos la mitad de los alumnos ha sacado un 7 o más.
  \task Las notas de la clase A son más parecidas entre sí.
  \task En la clase A hay alumnos con un 10.
\end{tasks}

\vfill
\begin{center}\color{muted}\small\faCheckCircle[regular]\ \ ¡Ánimo! Si te cuesta algún ejercicio, repasa los apuntes interactivos de la unidad.\end{center}

\end{document}
